\section{Dyadic geometry and the surviving right wing}
\label{sec:dyadic-corridor}

Choose a nonnegative $\psi\in C_c^\infty([1/2,2])$ satisfying
\begin{equation}
 \sum_{j\in\Z}\psi(x/2^j)=1
 \qquad(x>0).
 \label{eq:dyadic-partition}
\end{equation}
For dyadic $L,K$, put
\begin{align}
 \cH^\sharp_{L,K}(h;X)
 ={}&\sum_{\substack{\ell>M\\k>1}}
 \Lambda(\ell)r_R(\ell k+h)W(\ell k/X)
 \psi(\ell/L)\psi(k/K).
 \label{eq:dyadic-Hsharp-block}
\end{align}

\begin{proposition}[Exact dyadic corridor]
\label{prop:dyadic-corridor}
One has the finite exact decomposition
\begin{equation}
 \cH^\sharp_{h,R;M}(X)=
 \sum_{L,K}\cH^\sharp_{L,K}(h;X).
 \label{eq:exact-dyadic-decomposition}
\end{equation}
If $\supp W\subset[a,b]\subset(0,\infty)$, every nonzero block satisfies
\begin{equation}
 \frac{aX}{4}\le LK\le4bX,
 \qquad L>M/2,
 \qquad K\ll_W X/M.
 \label{eq:dyadic-support-ledger}
\end{equation}
There are $O_W(\log X)$ relevant pairs $(L,K)$.
\end{proposition}

\begin{proof}
Insert \eqref{eq:dyadic-partition} in both variables.  Compact support of
$W$ and the lower cutoffs make the sum finite.  If the two dyadic factors
are nonzero, then $\ell\asymp L$, $k\asymp K$, and
$\ell k\in[aX,bX]$, giving the product bounds.  For each $L$, the product
condition permits only $O_W(1)$ dyadic values of $K$.  There are
$O(\log X)$ possible $L$ between $M/2$ and $O_W(X)$.
\end{proof}

The normal form removes every physical term with $\ell\le M$.
Consequently every dyadic block with $L\le M/2$ vanishes, while every
surviving block satisfies $L>M/2$.  Up to the fixed support constants of
the dyadic partition, the closed left wing is therefore
\begin{equation}
 L\ll X^{1/2}\exp(-\sqrt{\log X}).
 \label{eq:closed-left-wing}
\end{equation}
The surviving set divides naturally into
\begin{align}
 &\text{near-square-root core:}
 &&X^{1/2}\e^{-\sqrt{\log X}}
 \ll L,K\ll
 X^{1/2}\e^{\sqrt{\log X}},
 \label{eq:near-square-root-core}\\
 &\text{right wing:}
 &&L\gg X^{1/2}\e^{\sqrt{\log X}},
 \qquad
 K\ll X^{1/2}\e^{-\sqrt{\log X}}.
 \label{eq:right-wing}
\end{align}
Here and below $\e^x$ denotes the ordinary exponential in prose-style
scale displays; it is unrelated to the additive character
$\e(t)=\exp(2\pi it)$ used later.

\begin{corollary}[Dyadic witness]
\label{cor:dyadic-witness}
If along a sequence $X_j$ one has
$|\cH^\sharp_{h,R;M}(X_j)|\ge\eta X_j$ for some $\eta>0$, then for every
such $X_j$ there is a relevant dyadic pair $(L_j,K_j)$ with
\begin{equation}
 |\cH^\sharp_{L,K}(h;X_j)|
 \gg_W\frac{\eta X_j}{\log X_j}.
 \label{eq:dyadic-witness-bound}
\end{equation}
Conversely, the uniform estimate
\begin{equation}
 \max_{L,K}|\cH^\sharp_{L,K}(h;X)|
 =o\left(\frac X{\log X}\right)
 \label{eq:uniform-block-sufficient}
\end{equation}
implies $\cH^\sharp=o(X)$.
\end{corollary}

The first assertion is only a pigeonhole statement.  It does not identify
one fixed dyadic scale across all $X$, and the converse condition is
sufficient rather than necessary because blocks can cancel.

\subsection{A fixed quotient is already a two-form problem}

It is tempting to call \eqref{eq:right-wing} easy because $k$ is small.
The coefficient $\Lambda(\ell)$, however, depends on the quotient in the
row $k\ell+h$.  Even a fixed slice contains an unresolved two-linear-form
correlation.

Fix $k\ge2$ and put $Y=X/k$.  Define
\begin{align}
 \cC_{k,h}(Y;W)
 &=\sum_{\ell\ge1}\Lambda(\ell)\Lambda(k\ell+h)W(\ell/Y),
 \label{eq:two-form-correlation}\\
 \cE_{k,h;R}(Y;W)
 &=\sum_{\ell\ge1}\Lambda(\ell)r_R(k\ell+h)W(\ell/Y).
 \label{eq:fixed-k-residual-slice}
\end{align}
Let
\begin{equation}
 \mathfrak S(k,h)=
 \prod_p\frac{1-\nu_p(k,h)/p}{(1-1/p)^2},
 \quad
 \nu_p(k,h)=
 \#\{a\pmod p:a(ka+h)\equiv0\pmod p\}.
 \label{eq:two-form-singular-series}
\end{equation}
For the finite model, put
\begin{equation}
 S_q(k,h)=
 \sum_{\substack{a\ ({\rm mod}\ q)\\(a,q)=1}}c_q(ka+h),
 \qquad
 \mathfrak S_R(k,h)=
 \sum_{q\le R}\frac{\mu(q)S_q(k,h)}{\varphi(q)^2}.
 \label{eq:finite-two-form-singular-series}
\end{equation}

\begin{proposition}[Fixed-$k$ boundary]
\label{prop:fixed-k-boundary}
Assume that the two forms $n$ and $kn+h$ are admissible.  For fixed
$k,h$ and $R=X^{1/2-\delta}$,
\begin{equation}
 \sum_\ell\Lambda(\ell)\Lambda_R(k\ell+h)W(\ell/Y)
 =\mathfrak S(k,h)YI_0(W)+o_{k,h,W}(Y).
 \label{eq:fixed-k-model-cross}
\end{equation}
Consequently
\begin{equation}
 \cC_{k,h}(Y;W)\sim\mathfrak S(k,h)YI_0(W)
 \quad\Longleftrightarrow\quad
 \cE_{k,h;R}(Y;W)=o(Y),
 \label{eq:fixed-k-equivalence}
\end{equation}
provided $I_0(W)\ne0$.
\end{proposition}

\begin{proof}
Admissibility implies $(k,h)=1$.  Since $k$ is fixed and $X=kY$,
\begin{equation}
 R=(kY)^{1/2-\delta}+O(1)
 \le\frac{Y^{1/2}}{(\log Y)^B}
 \label{eq:fixed-k-R-range}
\end{equation}
for every fixed $B$ and all sufficiently large $Y$.

Insert the divisor expansion
\eqref{eq:model-divisor-expansion-short}.  If $(d,k)>1$, the congruence
$d\mid k\ell+h$ has no solution.  If $(d,k)=1$, it is the single class
\begin{equation}
 \ell\equiv-h\overline{k}\pmod d,
 \label{eq:fixed-k-residue-class}
\end{equation}
which is reduced precisely when $(d,h)=1$.  Maximal
Bombieri--Vinogradov, \eqref{eq:fixed-k-R-range}, and
\eqref{eq:model-divisor-coefficient-pointwise} therefore give
\begin{align}
 &\sum_\ell\Lambda(\ell)\Lambda_R(k\ell+h)W(\ell/Y)
 \notag\\
 &\qquad={}
 YI_0(W)
 \sum_{\substack{d\le R\\(d,kh)=1}}
 \frac{\lambda'_R(d)}{\varphi(d)}
 +O_{A,k,h,W}\bigl(Y(\log Y)^{-A}\bigr).
 \label{eq:fixed-k-divisor-main}
\end{align}
Indeed, the coefficient costs only two additional logarithmic powers.
If $(d,k)=1$ but $(d,h)>1$, a nonzero source weight forces
$\ell=p^\nu$ for one of the finitely many $p\mid h$.  Only
$O_{h,W}(1)$ exponents lie on the support, and summing the divisor
coefficients over divisors of $kp^\nu+h$ contributes
$O_{k,h,W,\eps}(Y^\eps)$.

A finite rearrangement of the divisor and Ramanujan expansions gives
\begin{equation}
 \sum_{\substack{d\le R\\(d,kh)=1}}
 \frac{\lambda'_R(d)}{\varphi(d)}
 =\mathfrak S_R(k,h).
 \label{eq:fixed-k-finite-rearrangement}
\end{equation}
For completeness, this is also obtained by averaging
$\Lambda_R(ka+h)$ over reduced residue classes $a$ at each Ramanujan
denominator.

For squarefree $q$, the summand in
\eqref{eq:finite-two-form-singular-series} is multiplicative.  Its prime
values are
\begin{equation}
 \frac{\mu(p)S_p(k,h)}{\varphi(p)^2}
 =\begin{cases}
 -1/(p-1)^2,&p\nmid kh,\\
 1/(p-1),&p\mid k,\ p\nmid h,\\
 1/(p-1),&p\mid h,\ p\nmid k,\\
 -1,&p\mid(k,h).
 \end{cases}
 \label{eq:fixed-k-local-Ramanujan-values}
\end{equation}
Thus one plus the local value is exactly
$(1-\nu_p(k,h)/p)/(1-1/p)^2$.  The last case is absent under
admissibility.  The infinite Ramanujan series is absolutely convergent,
equals \eqref{eq:two-form-singular-series}, and satisfies
\begin{equation}
 \mathfrak S_R(k,h)=\mathfrak S(k,h)
 +O_{k,h}\left(\frac{(\log(2R))^2}{R}\right).
 \label{eq:fixed-k-singular-tail}
\end{equation}
Combining \eqref{eq:fixed-k-divisor-main}--
\eqref{eq:fixed-k-singular-tail} proves
\eqref{eq:fixed-k-model-cross}.  Finally, split
$\Lambda(k\ell+h)=\Lambda_R(k\ell+h)+r_R(k\ell+h)$ to obtain
\eqref{eq:fixed-k-equivalence}.
\end{proof}

For example, when $h=2$ and $k$ is odd, the forms $n$ and $kn+2$ are
admissible.  Their conjectural simultaneous primality is a special case of
the Hardy--Littlewood/Dickson prime-tuples conjecture
\citep{HardyLittlewood1923}.  The proposition does not say that a bound
for the total right wing must solve every fixed-$k$ slice separately;
cross-$k$ cancellation remains possible.  It does prove that the right
wing cannot be discarded merely by relabeling it Type I.
